
\documentclass{article}
\usepackage{colortbl}
\usepackage{makecell}
\usepackage{multirow}
\usepackage{supertabular}

\begin{document}

\newcounter{utterance}

\centering \large Interaction Transcript for game `clean\_up', experiment `2\_hard\_7obj\_et', episode 3 with glm{-}5.2.
\vspace{24pt}

{ \footnotesize  \setcounter{utterance}{1}
\setlength{\tabcolsep}{0pt}
\begin{supertabular}{c@{$\;$}|p{.15\linewidth}@{}p{.15\linewidth}p{.15\linewidth}p{.15\linewidth}p{.15\linewidth}p{.15\linewidth}}
    \# & $\;$A & \multicolumn{4}{c}{Game Master} & $\;\:$B\\
    \hline

    \theutterance \stepcounter{utterance}  
    & & \multicolumn{4}{p{0.6\linewidth}}{
        \cellcolor[rgb]{0.9,0.9,0.9}{
            \makecell[{{p{\linewidth}}}]{
                \texttt{\tiny{[P1$\langle$GM]}}
                \texttt{Ma olen teie mängujuht ja te mängite ühismängu järgmise ruudustikuga mängulauana:} \\
\\ 
\texttt{```} \\
\texttt{ 1234567} \\
\texttt{╔═╤═╤═══╗} \\
\texttt{║◌│◌│◌◌◌║ 1} \\
\texttt{║S├─┤T◌D║ 2} \\
\texttt{╟─┤◌├─┬─╢ 3} \\
\texttt{║◌│M│O│A║ 4} \\
\texttt{╟─┼─┘P│◌║ 5} \\
\texttt{║◌│◌◌┌┴─╢ 6} \\
\texttt{║◌│◌◌│◌◌║ 7} \\
\texttt{╚═╧══╧══╝} \\
\\ 
\texttt{```} \\
\\ 
\texttt{* Ülemisel serval on x{-}koordinaadid, mis suurenevad paremale, ja paremal serval y{-}koordinaadid, mis suurenevad allapoole.} \\
\texttt{* Järgmised objektid on teie ruudustikule juhuslikult paigutatud: 'P', 'O', 'T', 'S', 'D', 'A', 'M'.} \\
\\ 
\\ 
\texttt{Teine mängija näeb mängulaua varianti, kus objektid on paigutatud erinevatesse juhuslikesse asukohtadesse. Te ei näe teise mängija mängulauda ja tema ei näe teie mängulauda.} \\
\\ 
\\ 
\texttt{**Eesmärk:**} \\
\\ 
\texttt{Mõlemad mängijad peavad oma taustal olevaid objekte liigutama nii, et identsed objektid satuksid samadesse koordinaatidesse. Te peate teise mängijaga suhtlema, et kokku leppida ühises eesmärgi konfiguratsioonis.} \\
\\ 
\\ 
\texttt{**Reeglid:**} \\
\\ 
\texttt{* Igal käigul saate saata täpselt ühe kahest järgmisest käsust:} \\
\texttt{ 1. `ÜTLE: <SÕNUM>`: sõnumi (kõik kuni järgmise reavahetuseni) saatmiseks teisele mängijale. Ma edastan selle teie partnerile.} \\
\texttt{2. `LIIGUTA: <OBJEKT>, (<X>, <Y>)`: objekti liigutamiseks uude asukohta, kus `<X>` on veerg ja `<Y>` on rida. Ma teatan teile, kas teie käik õnnestus või mitte.} \\
\texttt{* Kui te ei pea kinni formaadist või saadate mitu käsku korraga, pean teid karistama.} \\
\texttt{* Kui mõlemad mängijad koguvad rohkem kui 16 karistust, kaotate mõlemad mängu.} \\
\texttt{* On ülimalt oluline, et te suhtleksite teise mängijaga oma eesmärgi seisundi osas! *Ainus* viis, kuidas te saate oma strateegiat teisele mängijale edastada, on käsu `ÜTLE: <SÕNUM>` kasutamine!} \\
\\ 
\\ 
\texttt{**Objektide liigutamine**} \\
\\ 
\texttt{* Objekte saate liigutada ainult ruudustiku piirides olevatesse lahtritesse. Sihtlahter peab olema tühi, st see peab sisaldama ainult sümbolit '◌'.} \\
\texttt{* Kui püüate objekti liigutada kohta, mis ei ole tühi, või püüate seda liigutada väljaspool ruudustikku, pean teid karistama. Saate uue katse.} \\
\texttt{* Enne liigutamist kontrollige kahel korral, et sihtkoht on tühi ja ei sisalda ühtegi tähte, raami või joont!} \\
\\ 
\\ 
\texttt{**Mängu lõpp**} \\
\\ 
\texttt{Kui arvate, et olete saavutanud kõigi objektide joondamise eesmärgi, võite paluda teisel mängijal mäng lõpetada, saates `ÜTLE: valmis?`. Kui teine mängija palub teil mängu lõpetada ja te vastate `ÜTLE: valmis!`, siis mäng lõpeb.} \\
\\ 
\texttt{Mõlemad mängijad võidavad, kui mäng lõpeb 28 vooru jooksul, kus üks voor on määratletud kahe mängija poolt ühe kehtiva käsu saatmisena.} \\
\\ 
\texttt{**Punktiarvestus:**} \\
\\ 
\texttt{Mida lähemal on identsed objektid mõlemal mängulaual, seda rohkem punkte saate. Karistused vähendavad teie punkte. Kas suudate rekordit ületada?} \\
\\ 
\\ 
\texttt{Palun saada teisele mängijale sõnum, et mängu alustada!} \\
            }
        }
    }
    & \\ \\

    \theutterance \stepcounter{utterance}  
    & \multicolumn{4}{p{0.6\linewidth}}{
        \cellcolor[rgb]{0.9,0.9,0.9}{
            \makecell[{{p{\linewidth}}}]{
                \texttt{\tiny{[P1$\rangle$GM]}}
                \texttt{ÜTLE: Tere! Alustame koostööd. Minu mängulaual on 7 objekti: P, O, T, S, D, A, M. Eesmärk on paigutada need samadele koordinaatidele. Minu praegused asukohad: S(1,2), T(3,2), D(7,2), M(2,4), O(3,4), A(5,4), P(5,5). Palun saadetage oma objektide asukohad, et saaksime leida ühise lahenduse. Üks lihtne variant oleks paigutada kõik tühjadesse lahtritesse ülemisse vasakusse ossa, näiteks: S(1,1), T(2,1), D(3,1), M(4,1), O(5,1), A(6,1), P(7,1). Andke teada, kas see sobib või milline on teie idee.} \\
            }
        }
    }
    & & \\ \\

    \theutterance \stepcounter{utterance}  
    & & & \multicolumn{2}{p{0.3\linewidth}}{
        \cellcolor[rgb]{0.9,0.9,0.9}{
            \makecell[{{p{\linewidth}}}]{
                \texttt{\tiny{[GM$|$GM]}}
                \texttt{ÜTLE: Tere! Alustame koostööd. Minu mängulaual on 7 objekti: P, O, T, S, D, A, M. Eesmärk on paigutada need samadele koordinaatidele. Minu praegused asukohad: S(1,2), T(3,2), D(7,2), M(2,4), O(3,4), A(5,4), P(5,5). Palun saadetage oma objektide asukohad, et saaksime leida ühise lahenduse. Üks lihtne variant oleks paigutada kõik tühjadesse lahtritesse ülemisse vasakusse ossa, näiteks: S(1,1), T(2,1), D(3,1), M(4,1), O(5,1), A(6,1), P(7,1). Andke teada, kas see sobib või milline on teie idee.} \\
            }
        }
    }
    & & \\ \\

    \theutterance \stepcounter{utterance}  
    & & \multicolumn{4}{p{0.6\linewidth}}{
        \cellcolor[rgb]{0.9,0.9,0.9}{
            \makecell[{{p{\linewidth}}}]{
                \texttt{\tiny{[GM$\rangle$P2]}}
                \texttt{Ma olen teie mängujuht ja te mängite ühismängu järgmise ruudustikuga mängulauana:} \\
\\ 
\texttt{```} \\
\texttt{ 1234567} \\
\texttt{╔═╤═╤═══╗} \\
\texttt{║◌│◌│◌PO║ 1} \\
\texttt{║◌├─┤◌◌◌║ 2} \\
\texttt{╟─┤◌├─┬─╢ 3} \\
\texttt{║◌│◌│◌│A║ 4} \\
\texttt{╟─┼─┘D│◌║ 5} \\
\texttt{║T│◌◌┌┴─╢ 6} \\
\texttt{║M│◌◌│◌S║ 7} \\
\texttt{╚═╧══╧══╝} \\
\\ 
\texttt{```} \\
\\ 
\texttt{* Ülemisel serval on x{-}koordinaadid, mis suurenevad paremale, ja paremal serval y{-}koordinaadid, mis suurenevad allapoole.} \\
\texttt{* Järgmised objektid on teie ruudustikule juhuslikult paigutatud: 'P', 'O', 'T', 'S', 'D', 'A', 'M'.} \\
\\ 
\\ 
\texttt{Teine mängija näeb mängulaua varianti, kus objektid on paigutatud erinevatesse juhuslikesse asukohtadesse. Te ei näe teise mängija mängulauda ja tema ei näe teie mängulauda.} \\
\\ 
\\ 
\texttt{**Eesmärk:**} \\
\\ 
\texttt{Mõlemad mängijad peavad oma taustal olevaid objekte liigutama nii, et identsed objektid satuksid samadesse koordinaatidesse. Te peate teise mängijaga suhtlema, et kokku leppida ühises eesmärgi konfiguratsioonis.} \\
\\ 
\\ 
\texttt{**Reeglid:**} \\
\\ 
\texttt{* Igal käigul saate saata täpselt ühe kahest järgmisest käsust:} \\
\texttt{ 1. `ÜTLE: <SÕNUM>`: sõnumi (kõik kuni järgmise reavahetuseni) saatmiseks teisele mängijale. Ma edastan selle teie partnerile.} \\
\texttt{2. `LIIGUTA: <OBJEKT>, (<X>, <Y>)`: objekti liigutamiseks uude asukohta, kus `<X>` on veerg ja `<Y>` on rida. Ma teatan teile, kas teie käik õnnestus või mitte.} \\
\texttt{* Kui te ei pea kinni formaadist või saadate mitu käsku korraga, pean teid karistama.} \\
\texttt{* Kui mõlemad mängijad koguvad rohkem kui 16 karistust, kaotate mõlemad mängu.} \\
\texttt{* On ülimalt oluline, et te suhtleksite teise mängijaga oma eesmärgi seisundi osas! *Ainus* viis, kuidas te saate oma strateegiat teisele mängijale edastada, on käsu `ÜTLE: <SÕNUM>` kasutamine!} \\
\\ 
\\ 
\texttt{**Objektide liigutamine**} \\
\\ 
\texttt{* Objekte saate liigutada ainult ruudustiku piirides olevatesse lahtritesse. Sihtlahter peab olema tühi, st see peab sisaldama ainult sümbolit '◌'.} \\
\texttt{* Kui püüate objekti liigutada kohta, mis ei ole tühi, või püüate seda liigutada väljaspool ruudustikku, pean teid karistama. Saate uue katse.} \\
\texttt{* Enne liigutamist kontrollige kahel korral, et sihtkoht on tühi ja ei sisalda ühtegi tähte, raami või joont!} \\
\\ 
\\ 
\texttt{**Mängu lõpp**} \\
\\ 
\texttt{Kui arvate, et olete saavutanud kõigi objektide joondamise eesmärgi, võite paluda teisel mängijal mäng lõpetada, saates `ÜTLE: valmis?`. Kui teine mängija palub teil mängu lõpetada ja te vastate `ÜTLE: valmis!`, siis mäng lõpeb.} \\
\\ 
\texttt{Mõlemad mängijad võidavad, kui mäng lõpeb 28 vooru jooksul, kus üks voor on määratletud kahe mängija poolt ühe kehtiva käsu saatmisena.} \\
\\ 
\texttt{**Punktiarvestus:**} \\
\\ 
\texttt{Mida lähemal on identsed objektid mõlemal mängulaual, seda rohkem punkte saate. Karistused vähendavad teie punkte. Kas suudate rekordit ületada?} \\
\\ 
\texttt{Teine mängija alustas mängu selle sõnumi saatmisega:} \\
\\ 
\texttt{"Tere! Alustame koostööd. Minu mängulaual on 7 objekti: P, O, T, S, D, A, M. Eesmärk on paigutada need samadele koordinaatidele. Minu praegused asukohad: S(1,2), T(3,2), D(7,2), M(2,4), O(3,4), A(5,4), P(5,5). Palun saadetage oma objektide asukohad, et saaksime leida ühise lahenduse. Üks lihtne variant oleks paigutada kõik tühjadesse lahtritesse ülemisse vasakusse ossa, näiteks: S(1,1), T(2,1), D(3,1), M(4,1), O(5,1), A(6,1), P(7,1). Andke teada, kas see sobib või milline on teie idee."} \\
\\ 
\texttt{Mis on sinu esimene käsk?} \\
            }
        }
    }
    & \\ \\

    \theutterance \stepcounter{utterance}  
    & & & \multicolumn{4}{p{0.6\linewidth}}{
        \cellcolor[rgb]{0.9,0.9,0.9}{
            \makecell[{{p{\linewidth}}}]{
                \texttt{\tiny{[GM$\langle$P2]}}
                \texttt{ÜTLE: Tere! Ait Lahendus sobib. Panen objektid: P(1,1), O(2,1), T(3,1), S(4,1), D(5,1), A(6,1), M(7,1). Alustan liigutamisega. Palun alusta ka sina.} \\
            }
        }
    }
     \\ \\

    \theutterance \stepcounter{utterance}  
    & & & \multicolumn{2}{p{0.3\linewidth}}{
        \cellcolor[rgb]{0.9,0.9,0.9}{
            \makecell[{{p{\linewidth}}}]{
                \texttt{\tiny{[GM$|$GM]}}
                \texttt{ÜTLE: Tere! Ait Lahendus sobib. Panen objektid: P(1,1), O(2,1), T(3,1), S(4,1), D(5,1), A(6,1), M(7,1). Alustan liigutamisega. Palun alusta ka sina.} \\
            }
        }
    }
    & & \\ \\

    \theutterance \stepcounter{utterance}  
    & & \multicolumn{4}{p{0.6\linewidth}}{
        \cellcolor[rgb]{0.9,0.9,0.9}{
            \makecell[{{p{\linewidth}}}]{
                \texttt{\tiny{[P1$\langle$GM]}}
                \texttt{Sinu sõnum on teisele mängijale edastatud.} \\
\\ 
\\ 
\texttt{Sa mängid praegu vooru 1 maksimaalsest 28 voorust.} \\
\texttt{Teil on kollektiivselt kogunenud 0 karistust 16 karistusest.} \\
\texttt{Teine mängija saatis selle sõnumi:} \\
\\ 
\\ 
\texttt{"Tere! Ait Lahendus sobib. Panen objektid: P(1,1), O(2,1), T(3,1), S(4,1), D(5,1), A(6,1), M(7,1). Alustan liigutamisega. Palun alusta ka sina."} \\
\\ 
\\ 
\texttt{Mis on sinu järgmine käsk?} \\
            }
        }
    }
    & \\ \\

    \theutterance \stepcounter{utterance}  
    & \multicolumn{4}{p{0.6\linewidth}}{
        \cellcolor[rgb]{0.9,0.9,0.9}{
            \makecell[{{p{\linewidth}}}]{
                \texttt{\tiny{[P1$\rangle$GM]}}
                \texttt{ÜTLE: Tere! Aitäh, et nõustusid lahendusega. Ootan, et me mõlemad liigutame objektid samadele lõplikele asukohtadele: P(1,1), O(2,1), T(3,1), S(4,1), D(5,1), A(6,1), M(7,1). Alustan nüüd oma objektide liigutamisest.} \\
            }
        }
    }
    & & \\ \\

    \theutterance \stepcounter{utterance}  
    & & & \multicolumn{2}{p{0.3\linewidth}}{
        \cellcolor[rgb]{0.9,0.9,0.9}{
            \makecell[{{p{\linewidth}}}]{
                \texttt{\tiny{[GM$|$GM]}}
                \texttt{ÜTLE: Tere! Aitäh, et nõustusid lahendusega. Ootan, et me mõlemad liigutame objektid samadele lõplikele asukohtadele: P(1,1), O(2,1), T(3,1), S(4,1), D(5,1), A(6,1), M(7,1). Alustan nüüd oma objektide liigutamisest.} \\
            }
        }
    }
    & & \\ \\

    \theutterance \stepcounter{utterance}  
    & & \multicolumn{4}{p{0.6\linewidth}}{
        \cellcolor[rgb]{0.9,0.9,0.9}{
            \makecell[{{p{\linewidth}}}]{
                \texttt{\tiny{[GM$\rangle$P2]}}
                \texttt{Sinu sõnum on teisele mängijale edastatud.} \\
\\ 
\\ 
\texttt{Sa mängid praegu vooru 2 maksimaalsest 28 voorust.} \\
\texttt{Teil on kollektiivselt kogunenud 0 karistust 16 karistusest.} \\
\texttt{Teine mängija saatis selle sõnumi:} \\
\\ 
\\ 
\texttt{"Tere! Aitäh, et nõustusid lahendusega. Ootan, et me mõlemad liigutame objektid samadele lõplikele asukohtadele: P(1,1), O(2,1), T(3,1), S(4,1), D(5,1), A(6,1), M(7,1). Alustan nüüd oma objektide liigutamisest."} \\
\\ 
\\ 
\texttt{Mis on sinu järgmine käsk?} \\
            }
        }
    }
    & \\ \\

    \theutterance \stepcounter{utterance}  
    & & & \multicolumn{4}{p{0.6\linewidth}}{
        \cellcolor[rgb]{0.9,0.9,0.9}{
            \makecell[{{p{\linewidth}}}]{
                \texttt{\tiny{[GM$\langle$P2]}}
                \texttt{LIIGUTA: P, (1, 1)} \\
            }
        }
    }
     \\ \\

    \theutterance \stepcounter{utterance}  
    & & & \multicolumn{2}{p{0.3\linewidth}}{
        \cellcolor[rgb]{0.9,0.9,0.9}{
            \makecell[{{p{\linewidth}}}]{
                \texttt{\tiny{[GM$|$GM]}}
                \texttt{LIIGUTA: P, (1, 1)} \\
            }
        }
    }
    & & \\ \\

    \theutterance \stepcounter{utterance}  
    & & & \multicolumn{2}{p{0.3\linewidth}}{
        \cellcolor[rgb]{0.9,0.9,0.9}{
            \makecell[{{p{\linewidth}}}]{
                \texttt{\tiny{[GM$|$GM]}}
                \texttt{P moved to (1, 1)} \\
            }
        }
    }
    & & \\ \\

    \theutterance \stepcounter{utterance}  
    & & \multicolumn{4}{p{0.6\linewidth}}{
        \cellcolor[rgb]{0.9,0.9,0.9}{
            \makecell[{{p{\linewidth}}}]{
                \texttt{\tiny{[P1$\langle$GM]}}
                \texttt{Sinu sõnum on teisele mängijale edastatud.} \\
\\ 
\\ 
\texttt{Sa mängid praegu vooru 2 maksimaalsest 28 voorust.} \\
\texttt{Teil on kollektiivselt kogunenud 0 karistust 16 karistusest.} \\
\texttt{Teine mängija liigutas objekti oma ruudustikul.} \\
\\ 
\\ 
\texttt{Mis on sinu järgmine käsk?} \\
            }
        }
    }
    & \\ \\

    \theutterance \stepcounter{utterance}  
    & \multicolumn{4}{p{0.6\linewidth}}{
        \cellcolor[rgb]{0.9,0.9,0.9}{
            \makecell[{{p{\linewidth}}}]{
                \texttt{\tiny{[P1$\rangle$GM]}}
                \texttt{LIIGUTA: S, (1, 1)} \\
            }
        }
    }
    & & \\ \\

    \theutterance \stepcounter{utterance}  
    & & & \multicolumn{2}{p{0.3\linewidth}}{
        \cellcolor[rgb]{0.9,0.9,0.9}{
            \makecell[{{p{\linewidth}}}]{
                \texttt{\tiny{[GM$|$GM]}}
                \texttt{LIIGUTA: S, (1, 1)} \\
            }
        }
    }
    & & \\ \\

    \theutterance \stepcounter{utterance}  
    & & & \multicolumn{2}{p{0.3\linewidth}}{
        \cellcolor[rgb]{0.9,0.9,0.9}{
            \makecell[{{p{\linewidth}}}]{
                \texttt{\tiny{[GM$|$GM]}}
                \texttt{S moved to (1, 1)} \\
            }
        }
    }
    & & \\ \\

    \theutterance \stepcounter{utterance}  
    & & \multicolumn{4}{p{0.6\linewidth}}{
        \cellcolor[rgb]{0.9,0.9,0.9}{
            \makecell[{{p{\linewidth}}}]{
                \texttt{\tiny{[GM$\rangle$P2]}}
                \texttt{Edukalt teisaldatud 'P' asukohta (1,1). Teie uuendatud ruudustik näeb välja selline:} \\
\texttt{```} \\
\texttt{ 1234567} \\
\texttt{╔═╤═╤═══╗} \\
\texttt{║P│◌│◌◌O║ 1} \\
\texttt{║◌├─┤◌◌◌║ 2} \\
\texttt{╟─┤◌├─┬─╢ 3} \\
\texttt{║◌│◌│◌│A║ 4} \\
\texttt{╟─┼─┘D│◌║ 5} \\
\texttt{║T│◌◌┌┴─╢ 6} \\
\texttt{║M│◌◌│◌S║ 7} \\
\texttt{╚═╧══╧══╝} \\
\texttt{```} \\
\texttt{Sa mängid praegu vooru 3 maksimaalsest 28 voorust.} \\
\texttt{Teil on kollektiivselt kogunenud 0 karistust 16 karistusest.} \\
\texttt{Teine mängija liigutas objekti oma ruudustikul.} \\
\\ 
\\ 
\texttt{Mis on sinu järgmine käsk?} \\
            }
        }
    }
    & \\ \\

    \theutterance \stepcounter{utterance}  
    & & & \multicolumn{4}{p{0.6\linewidth}}{
        \cellcolor[rgb]{0.9,0.9,0.9}{
            \makecell[{{p{\linewidth}}}]{
                \texttt{\tiny{[GM$\langle$P2]}}
                \texttt{LIIGUTA: O, (2, 1)} \\
            }
        }
    }
     \\ \\

    \theutterance \stepcounter{utterance}  
    & & & \multicolumn{2}{p{0.3\linewidth}}{
        \cellcolor[rgb]{0.9,0.9,0.9}{
            \makecell[{{p{\linewidth}}}]{
                \texttt{\tiny{[GM$|$GM]}}
                \texttt{LIIGUTA: O, (2, 1)} \\
            }
        }
    }
    & & \\ \\

    \theutterance \stepcounter{utterance}  
    & & & \multicolumn{2}{p{0.3\linewidth}}{
        \cellcolor[rgb]{0.9,0.9,0.9}{
            \makecell[{{p{\linewidth}}}]{
                \texttt{\tiny{[GM$|$GM]}}
                \texttt{Karistus: (2,1) ei ole tühi, vaid sisaldab objekti '│'.} \\
\\ 
\texttt{Teil on kollektiivselt kogunenud 1 karistust 16 karistusest.} \\
\texttt{Palun proovi uuesti!} \\
            }
        }
    }
    & & \\ \\

    \theutterance \stepcounter{utterance}  
    & & \multicolumn{4}{p{0.6\linewidth}}{
        \cellcolor[rgb]{0.9,0.9,0.9}{
            \makecell[{{p{\linewidth}}}]{
                \texttt{\tiny{[GM$\rangle$P2]}}
                \texttt{Karistus: (2,1) ei ole tühi, vaid sisaldab objekti '│'.} \\
\\ 
\texttt{Teil on kollektiivselt kogunenud 1 karistust 16 karistusest.} \\
\texttt{Palun proovi uuesti!} \\
            }
        }
    }
    & \\ \\

    \theutterance \stepcounter{utterance}  
    & & & \multicolumn{4}{p{0.6\linewidth}}{
        \cellcolor[rgb]{0.9,0.9,0.9}{
            \makecell[{{p{\linewidth}}}]{
                \texttt{\tiny{[GM$\langle$P2]}}
                \texttt{Vabandan! Vaatasin ruudustikku hoolitsematult {-} (2,1) on joon │, mitte tühi lahter.} \\
\\ 
\texttt{Uurisin ruudustikku hooli} \\
            }
        }
    }
     \\ \\

    \theutterance \stepcounter{utterance}  
    & & & \multicolumn{2}{p{0.3\linewidth}}{
        \cellcolor[rgb]{0.9,0.9,0.9}{
            \makecell[{{p{\linewidth}}}]{
                \texttt{\tiny{[GM$|$GM]}}
                \texttt{Vabandan! Vaatasin ruudustikku hoolitsematult {-} (2,1) on joon │, mitte tühi lahter.} \\
\\ 
\texttt{Uurisin ruudustikku hooli} \\
            }
        }
    }
    & & \\ \\

    \theutterance \stepcounter{utterance}  
    & & & \multicolumn{2}{p{0.3\linewidth}}{
        \cellcolor[rgb]{0.9,0.9,0.9}{
            \makecell[{{p{\linewidth}}}]{
                \texttt{\tiny{[GM$|$GM]}}
                \texttt{Invalid response format} \\
            }
        }
    }
    & & \\ \\

    \theutterance \stepcounter{utterance}  
    & & \multicolumn{4}{p{0.6\linewidth}}{
        \cellcolor[rgb]{0.9,0.9,0.9}{
            \makecell[{{p{\linewidth}}}]{
                \texttt{\tiny{[GM$\rangle$P2]}}
                \texttt{Karistus: Teie sõnum ei ole oodatud formaadis!} \\
\\ 
\texttt{Veendu, et sinu vastus sisaldab ainult kas `ÜTLE: <SÕNUM>` või `LIIGUTA: <OBJEKT>, (<X>, <Y>)` ja mitte midagi muud!} \\
\\ 
\texttt{Teil on kollektiivselt kogunenud 2 karistust 16 karistusest.} \\
\texttt{Palun proovi uuesti!} \\
            }
        }
    }
    & \\ \\

    \theutterance \stepcounter{utterance}  
    & & & \multicolumn{4}{p{0.6\linewidth}}{
        \cellcolor[rgb]{0.9,0.9,0.9}{
            \makecell[{{p{\linewidth}}}]{
                \texttt{\tiny{[GM$\langle$P2]}}
                \texttt{LIIGUTA: O, (3, 1)} \\
            }
        }
    }
     \\ \\

    \theutterance \stepcounter{utterance}  
    & & & \multicolumn{2}{p{0.3\linewidth}}{
        \cellcolor[rgb]{0.9,0.9,0.9}{
            \makecell[{{p{\linewidth}}}]{
                \texttt{\tiny{[GM$|$GM]}}
                \texttt{LIIGUTA: O, (3, 1)} \\
            }
        }
    }
    & & \\ \\

    \theutterance \stepcounter{utterance}  
    & & & \multicolumn{2}{p{0.3\linewidth}}{
        \cellcolor[rgb]{0.9,0.9,0.9}{
            \makecell[{{p{\linewidth}}}]{
                \texttt{\tiny{[GM$|$GM]}}
                \texttt{O moved to (3, 1)} \\
            }
        }
    }
    & & \\ \\

    \theutterance \stepcounter{utterance}  
    & & \multicolumn{4}{p{0.6\linewidth}}{
        \cellcolor[rgb]{0.9,0.9,0.9}{
            \makecell[{{p{\linewidth}}}]{
                \texttt{\tiny{[P1$\langle$GM]}}
                \texttt{Edukalt teisaldatud 'S' asukohta (1,1). Teie uuendatud ruudustik näeb välja selline:} \\
\texttt{```} \\
\texttt{ 1234567} \\
\texttt{╔═╤═╤═══╗} \\
\texttt{║S│◌│◌◌◌║ 1} \\
\texttt{║◌├─┤T◌D║ 2} \\
\texttt{╟─┤◌├─┬─╢ 3} \\
\texttt{║◌│M│O│A║ 4} \\
\texttt{╟─┼─┘P│◌║ 5} \\
\texttt{║◌│◌◌┌┴─╢ 6} \\
\texttt{║◌│◌◌│◌◌║ 7} \\
\texttt{╚═╧══╧══╝} \\
\texttt{```} \\
\texttt{Sa mängid praegu vooru 3 maksimaalsest 28 voorust.} \\
\texttt{Teil on kollektiivselt kogunenud 2 karistust 16 karistusest.} \\
\texttt{Teine mängija liigutas objekti oma ruudustikul.} \\
\\ 
\\ 
\texttt{Mis on sinu järgmine käsk?} \\
            }
        }
    }
    & \\ \\

    \theutterance \stepcounter{utterance}  
    & \multicolumn{4}{p{0.6\linewidth}}{
        \cellcolor[rgb]{0.9,0.9,0.9}{
            \makecell[{{p{\linewidth}}}]{
                \texttt{\tiny{[P1$\rangle$GM]}}
                \texttt{LIIGUTA: P, (2, 1)} \\
            }
        }
    }
    & & \\ \\

    \theutterance \stepcounter{utterance}  
    & & & \multicolumn{2}{p{0.3\linewidth}}{
        \cellcolor[rgb]{0.9,0.9,0.9}{
            \makecell[{{p{\linewidth}}}]{
                \texttt{\tiny{[GM$|$GM]}}
                \texttt{LIIGUTA: P, (2, 1)} \\
            }
        }
    }
    & & \\ \\

    \theutterance \stepcounter{utterance}  
    & & & \multicolumn{2}{p{0.3\linewidth}}{
        \cellcolor[rgb]{0.9,0.9,0.9}{
            \makecell[{{p{\linewidth}}}]{
                \texttt{\tiny{[GM$|$GM]}}
                \texttt{Karistus: (2,1) ei ole tühi, vaid sisaldab objekti '│'.} \\
\\ 
\texttt{Teil on kollektiivselt kogunenud 3 karistust 16 karistusest.} \\
\texttt{Palun proovi uuesti!} \\
            }
        }
    }
    & & \\ \\

    \theutterance \stepcounter{utterance}  
    & & \multicolumn{4}{p{0.6\linewidth}}{
        \cellcolor[rgb]{0.9,0.9,0.9}{
            \makecell[{{p{\linewidth}}}]{
                \texttt{\tiny{[P1$\langle$GM]}}
                \texttt{Karistus: (2,1) ei ole tühi, vaid sisaldab objekti '│'.} \\
\\ 
\texttt{Teil on kollektiivselt kogunenud 3 karistust 16 karistusest.} \\
\texttt{Palun proovi uuesti!} \\
            }
        }
    }
    & \\ \\

    \theutterance \stepcounter{utterance}  
    & \multicolumn{4}{p{0.6\linewidth}}{
        \cellcolor[rgb]{0.9,0.9,0.9}{
            \makecell[{{p{\linewidth}}}]{
                \texttt{\tiny{[P1$\rangle$GM]}}
                \texttt{Tegin vea {-} (2,1) sisaldab raamisümbolit. Vaatan uuesti ruudustikku ja valin tühi lahtri.} \\
\\ 
\texttt{LIIGUTA: P, (4, 1)} \\
            }
        }
    }
    & & \\ \\

    \theutterance \stepcounter{utterance}  
    & & & \multicolumn{2}{p{0.3\linewidth}}{
        \cellcolor[rgb]{0.9,0.9,0.9}{
            \makecell[{{p{\linewidth}}}]{
                \texttt{\tiny{[GM$|$GM]}}
                \texttt{Tegin vea {-} (2,1) sisaldab raamisümbolit. Vaatan uuesti ruudustikku ja valin tühi lahtri.} \\
\\ 
\texttt{LIIGUTA: P, (4, 1)} \\
            }
        }
    }
    & & \\ \\

    \theutterance \stepcounter{utterance}  
    & & & \multicolumn{2}{p{0.3\linewidth}}{
        \cellcolor[rgb]{0.9,0.9,0.9}{
            \makecell[{{p{\linewidth}}}]{
                \texttt{\tiny{[GM$|$GM]}}
                \texttt{Invalid format: head is not empty: 'Tegin vea {-} (2,1) sisaldab raamisümbolit. Vaatan uuesti ruudustikku ja valin tühi lahtri.} \\
\\ 
\texttt{'} \\
            }
        }
    }
    & & \\ \\

    \theutterance \stepcounter{utterance}  
    & & \multicolumn{4}{p{0.6\linewidth}}{
        \cellcolor[rgb]{0.9,0.9,0.9}{
            \makecell[{{p{\linewidth}}}]{
                \texttt{\tiny{[P1$\langle$GM]}}
                \texttt{Karistus: Teie sõnum ei tohi sisaldada midagi enne käsku!} \\
\\ 
\texttt{Veendu, et sinu vastus sisaldab ainult kas `ÜTLE: <SÕNUM>` või `LIIGUTA: <OBJEKT>, (<X>, <Y>)` ja mitte midagi muud!} \\
\\ 
\texttt{Teil on kollektiivselt kogunenud 4 karistust 16 karistusest.} \\
\texttt{Palun proovi uuesti!} \\
            }
        }
    }
    & \\ \\

    \theutterance \stepcounter{utterance}  
    & \multicolumn{4}{p{0.6\linewidth}}{
        \cellcolor[rgb]{0.9,0.9,0.9}{
            \makecell[{{p{\linewidth}}}]{
                \texttt{\tiny{[P1$\rangle$GM]}}
                \texttt{LIIGUTA: P, (4, 1)} \\
            }
        }
    }
    & & \\ \\

    \theutterance \stepcounter{utterance}  
    & & & \multicolumn{2}{p{0.3\linewidth}}{
        \cellcolor[rgb]{0.9,0.9,0.9}{
            \makecell[{{p{\linewidth}}}]{
                \texttt{\tiny{[GM$|$GM]}}
                \texttt{LIIGUTA: P, (4, 1)} \\
            }
        }
    }
    & & \\ \\

    \theutterance \stepcounter{utterance}  
    & & & \multicolumn{2}{p{0.3\linewidth}}{
        \cellcolor[rgb]{0.9,0.9,0.9}{
            \makecell[{{p{\linewidth}}}]{
                \texttt{\tiny{[GM$|$GM]}}
                \texttt{Karistus: (4,1) ei ole tühi, vaid sisaldab objekti '│'.} \\
\\ 
\texttt{Teil on kollektiivselt kogunenud 5 karistust 16 karistusest.} \\
\texttt{Palun proovi uuesti!} \\
            }
        }
    }
    & & \\ \\

    \theutterance \stepcounter{utterance}  
    & & \multicolumn{4}{p{0.6\linewidth}}{
        \cellcolor[rgb]{0.9,0.9,0.9}{
            \makecell[{{p{\linewidth}}}]{
                \texttt{\tiny{[P1$\langle$GM]}}
                \texttt{Karistus: (4,1) ei ole tühi, vaid sisaldab objekti '│'.} \\
\\ 
\texttt{Teil on kollektiivselt kogunenud 5 karistust 16 karistusest.} \\
\texttt{Palun proovi uuesti!} \\
            }
        }
    }
    & \\ \\

    \theutterance \stepcounter{utterance}  
    & \multicolumn{4}{p{0.6\linewidth}}{
        \cellcolor[rgb]{0.9,0.9,0.9}{
            \makecell[{{p{\linewidth}}}]{
                \texttt{\tiny{[P1$\rangle$GM]}}
                \texttt{LIIGUTA: P, (7, 1)} \\
            }
        }
    }
    & & \\ \\

    \theutterance \stepcounter{utterance}  
    & & & \multicolumn{2}{p{0.3\linewidth}}{
        \cellcolor[rgb]{0.9,0.9,0.9}{
            \makecell[{{p{\linewidth}}}]{
                \texttt{\tiny{[GM$|$GM]}}
                \texttt{LIIGUTA: P, (7, 1)} \\
            }
        }
    }
    & & \\ \\

    \theutterance \stepcounter{utterance}  
    & & & \multicolumn{2}{p{0.3\linewidth}}{
        \cellcolor[rgb]{0.9,0.9,0.9}{
            \makecell[{{p{\linewidth}}}]{
                \texttt{\tiny{[GM$|$GM]}}
                \texttt{P moved to (7, 1)} \\
            }
        }
    }
    & & \\ \\

    \theutterance \stepcounter{utterance}  
    & & \multicolumn{4}{p{0.6\linewidth}}{
        \cellcolor[rgb]{0.9,0.9,0.9}{
            \makecell[{{p{\linewidth}}}]{
                \texttt{\tiny{[GM$\rangle$P2]}}
                \texttt{Edukalt teisaldatud 'O' asukohta (3,1). Teie uuendatud ruudustik näeb välja selline:} \\
\texttt{```} \\
\texttt{ 1234567} \\
\texttt{╔═╤═╤═══╗} \\
\texttt{║P│O│◌◌◌║ 1} \\
\texttt{║◌├─┤◌◌◌║ 2} \\
\texttt{╟─┤◌├─┬─╢ 3} \\
\texttt{║◌│◌│◌│A║ 4} \\
\texttt{╟─┼─┘D│◌║ 5} \\
\texttt{║T│◌◌┌┴─╢ 6} \\
\texttt{║M│◌◌│◌S║ 7} \\
\texttt{╚═╧══╧══╝} \\
\texttt{```} \\
\texttt{Sa mängid praegu vooru 4 maksimaalsest 28 voorust.} \\
\texttt{Teil on kollektiivselt kogunenud 5 karistust 16 karistusest.} \\
\texttt{Teine mängija liigutas objekti oma ruudustikul.} \\
\\ 
\\ 
\texttt{Mis on sinu järgmine käsk?} \\
            }
        }
    }
    & \\ \\

    \theutterance \stepcounter{utterance}  
    & & & \multicolumn{4}{p{0.6\linewidth}}{
        \cellcolor[rgb]{0.9,0.9,0.9}{
            \makecell[{{p{\linewidth}}}]{
                \texttt{\tiny{[GM$\langle$P2]}}
                \texttt{ÜTLE: Mu ruudustikul on (2,1) blokeeritud joonega, nii et O on nüüd (3,1). Vajan uut ühendkava. Pakun: P(1,1), T(4,1), O(5,1), S(6,1), D(4,2), A(5,2), M(4,3). Kas see sobib?} \\
            }
        }
    }
     \\ \\

    \theutterance \stepcounter{utterance}  
    & & & \multicolumn{2}{p{0.3\linewidth}}{
        \cellcolor[rgb]{0.9,0.9,0.9}{
            \makecell[{{p{\linewidth}}}]{
                \texttt{\tiny{[GM$|$GM]}}
                \texttt{ÜTLE: Mu ruudustikul on (2,1) blokeeritud joonega, nii et O on nüüd (3,1). Vajan uut ühendkava. Pakun: P(1,1), T(4,1), O(5,1), S(6,1), D(4,2), A(5,2), M(4,3). Kas see sobib?} \\
            }
        }
    }
    & & \\ \\

    \theutterance \stepcounter{utterance}  
    & & \multicolumn{4}{p{0.6\linewidth}}{
        \cellcolor[rgb]{0.9,0.9,0.9}{
            \makecell[{{p{\linewidth}}}]{
                \texttt{\tiny{[P1$\langle$GM]}}
                \texttt{Edukalt teisaldatud 'P' asukohta (7,1). Teie uuendatud ruudustik näeb välja selline:} \\
\texttt{```} \\
\texttt{ 1234567} \\
\texttt{╔═╤═╤═══╗} \\
\texttt{║S│◌│◌◌P║ 1} \\
\texttt{║◌├─┤T◌D║ 2} \\
\texttt{╟─┤◌├─┬─╢ 3} \\
\texttt{║◌│M│O│A║ 4} \\
\texttt{╟─┼─┘◌│◌║ 5} \\
\texttt{║◌│◌◌┌┴─╢ 6} \\
\texttt{║◌│◌◌│◌◌║ 7} \\
\texttt{╚═╧══╧══╝} \\
\texttt{```} \\
\texttt{Sa mängid praegu vooru 4 maksimaalsest 28 voorust.} \\
\texttt{Teil on kollektiivselt kogunenud 5 karistust 16 karistusest.} \\
\texttt{Teine mängija saatis selle sõnumi:} \\
\\ 
\\ 
\texttt{"Mu ruudustikul on (2,1) blokeeritud joonega, nii et O on nüüd (3,1). Vajan uut ühendkava. Pakun: P(1,1), T(4,1), O(5,1), S(6,1), D(4,2), A(5,2), M(4,3). Kas see sobib?"} \\
\\ 
\\ 
\texttt{Mis on sinu järgmine käsk?} \\
            }
        }
    }
    & \\ \\

    \theutterance \stepcounter{utterance}  
    & \multicolumn{4}{p{0.6\linewidth}}{
        \cellcolor[rgb]{0.9,0.9,0.9}{
            \makecell[{{p{\linewidth}}}]{
                \texttt{\tiny{[P1$\rangle$GM]}}
                \texttt{ÜTLE: Kahjuks on ka minu ruudustikul palju blokeeritud kohti. (2,1) ja (4,1) on blokeeritud joontega. Tühi lahtrid on: (3,1), (5,1), (6,1), (7,1), (1,2), (5,2), (6,2), (7,2), (1,3), (2,3), (1,4), (7,4), (4,5), (6,5), (1,6), (2,6), (3,6), (5,6), (1,7), (2,7), (3,7), (5,7), (6,7). Et vältida probleeme, paneme kõik ühte piirkonda: P(1,2), O(5,2), T(6,2), S(7,2), D(7,1), A(6,1), M(5,1). Kas see sobib?} \\
            }
        }
    }
    & & \\ \\

    \theutterance \stepcounter{utterance}  
    & & & \multicolumn{2}{p{0.3\linewidth}}{
        \cellcolor[rgb]{0.9,0.9,0.9}{
            \makecell[{{p{\linewidth}}}]{
                \texttt{\tiny{[GM$|$GM]}}
                \texttt{ÜTLE: Kahjuks on ka minu ruudustikul palju blokeeritud kohti. (2,1) ja (4,1) on blokeeritud joontega. Tühi lahtrid on: (3,1), (5,1), (6,1), (7,1), (1,2), (5,2), (6,2), (7,2), (1,3), (2,3), (1,4), (7,4), (4,5), (6,5), (1,6), (2,6), (3,6), (5,6), (1,7), (2,7), (3,7), (5,7), (6,7). Et vältida probleeme, paneme kõik ühte piirkonda: P(1,2), O(5,2), T(6,2), S(7,2), D(7,1), A(6,1), M(5,1). Kas see sobib?} \\
            }
        }
    }
    & & \\ \\

    \theutterance \stepcounter{utterance}  
    & & \multicolumn{4}{p{0.6\linewidth}}{
        \cellcolor[rgb]{0.9,0.9,0.9}{
            \makecell[{{p{\linewidth}}}]{
                \texttt{\tiny{[GM$\rangle$P2]}}
                \texttt{Sinu sõnum on teisele mängijale edastatud.} \\
\\ 
\\ 
\texttt{Sa mängid praegu vooru 5 maksimaalsest 28 voorust.} \\
\texttt{Teil on kollektiivselt kogunenud 5 karistust 16 karistusest.} \\
\texttt{Teine mängija saatis selle sõnumi:} \\
\\ 
\\ 
\texttt{"Kahjuks on ka minu ruudustikul palju blokeeritud kohti. (2,1) ja (4,1) on blokeeritud joontega. Tühi lahtrid on: (3,1), (5,1), (6,1), (7,1), (1,2), (5,2), (6,2), (7,2), (1,3), (2,3), (1,4), (7,4), (4,5), (6,5), (1,6), (2,6), (3,6), (5,6), (1,7), (2,7), (3,7), (5,7), (6,7). Et vältida probleeme, paneme kõik ühte piirkonda: P(1,2), O(5,2), T(6,2), S(7,2), D(7,1), A(6,1), M(5,1). Kas see sobib?"} \\
\\ 
\\ 
\texttt{Mis on sinu järgmine käsk?} \\
            }
        }
    }
    & \\ \\

    \theutterance \stepcounter{utterance}  
    & & & \multicolumn{4}{p{0.6\linewidth}}{
        \cellcolor[rgb]{0.9,0.9,0.9}{
            \makecell[{{p{\linewidth}}}]{
                \texttt{\tiny{[GM$\langle$P2]}}
                \texttt{ÜTLE: See sobib! Panen objektid: P(1,2), O(5,2), T(6,2), S(7,2), D(7,1), A(6,1), M(5,1). Alustan liigutamisest. Palun alusta ka sina.} \\
            }
        }
    }
     \\ \\

    \theutterance \stepcounter{utterance}  
    & & & \multicolumn{2}{p{0.3\linewidth}}{
        \cellcolor[rgb]{0.9,0.9,0.9}{
            \makecell[{{p{\linewidth}}}]{
                \texttt{\tiny{[GM$|$GM]}}
                \texttt{ÜTLE: See sobib! Panen objektid: P(1,2), O(5,2), T(6,2), S(7,2), D(7,1), A(6,1), M(5,1). Alustan liigutamisest. Palun alusta ka sina.} \\
            }
        }
    }
    & & \\ \\

    \theutterance \stepcounter{utterance}  
    & & \multicolumn{4}{p{0.6\linewidth}}{
        \cellcolor[rgb]{0.9,0.9,0.9}{
            \makecell[{{p{\linewidth}}}]{
                \texttt{\tiny{[P1$\langle$GM]}}
                \texttt{Sinu sõnum on teisele mängijale edastatud.} \\
\\ 
\\ 
\texttt{Sa mängid praegu vooru 5 maksimaalsest 28 voorust.} \\
\texttt{Teil on kollektiivselt kogunenud 5 karistust 16 karistusest.} \\
\texttt{Teine mängija saatis selle sõnumi:} \\
\\ 
\\ 
\texttt{"See sobib! Panen objektid: P(1,2), O(5,2), T(6,2), S(7,2), D(7,1), A(6,1), M(5,1). Alustan liigutamisest. Palun alusta ka sina."} \\
\\ 
\\ 
\texttt{Mis on sinu järgmine käsk?} \\
            }
        }
    }
    & \\ \\

    \theutterance \stepcounter{utterance}  
    & \multicolumn{4}{p{0.6\linewidth}}{
        \cellcolor[rgb]{0.9,0.9,0.9}{
            \makecell[{{p{\linewidth}}}]{
                \texttt{\tiny{[P1$\rangle$GM]}}
                \texttt{LIIGUTA: P, (1, 2)} \\
            }
        }
    }
    & & \\ \\

    \theutterance \stepcounter{utterance}  
    & & & \multicolumn{2}{p{0.3\linewidth}}{
        \cellcolor[rgb]{0.9,0.9,0.9}{
            \makecell[{{p{\linewidth}}}]{
                \texttt{\tiny{[GM$|$GM]}}
                \texttt{LIIGUTA: P, (1, 2)} \\
            }
        }
    }
    & & \\ \\

    \theutterance \stepcounter{utterance}  
    & & & \multicolumn{2}{p{0.3\linewidth}}{
        \cellcolor[rgb]{0.9,0.9,0.9}{
            \makecell[{{p{\linewidth}}}]{
                \texttt{\tiny{[GM$|$GM]}}
                \texttt{P moved to (1, 2)} \\
            }
        }
    }
    & & \\ \\

    \theutterance \stepcounter{utterance}  
    & & \multicolumn{4}{p{0.6\linewidth}}{
        \cellcolor[rgb]{0.9,0.9,0.9}{
            \makecell[{{p{\linewidth}}}]{
                \texttt{\tiny{[GM$\rangle$P2]}}
                \texttt{Sinu sõnum on teisele mängijale edastatud.} \\
\\ 
\\ 
\texttt{Sa mängid praegu vooru 6 maksimaalsest 28 voorust.} \\
\texttt{Teil on kollektiivselt kogunenud 5 karistust 16 karistusest.} \\
\texttt{Teine mängija liigutas objekti oma ruudustikul.} \\
\\ 
\\ 
\texttt{Mis on sinu järgmine käsk?} \\
            }
        }
    }
    & \\ \\

    \theutterance \stepcounter{utterance}  
    & & & \multicolumn{4}{p{0.6\linewidth}}{
        \cellcolor[rgb]{0.9,0.9,0.9}{
            \makecell[{{p{\linewidth}}}]{
                \texttt{\tiny{[GM$\langle$P2]}}
                \texttt{LIIGUTA: D, (7, 1)} \\
            }
        }
    }
     \\ \\

    \theutterance \stepcounter{utterance}  
    & & & \multicolumn{2}{p{0.3\linewidth}}{
        \cellcolor[rgb]{0.9,0.9,0.9}{
            \makecell[{{p{\linewidth}}}]{
                \texttt{\tiny{[GM$|$GM]}}
                \texttt{LIIGUTA: D, (7, 1)} \\
            }
        }
    }
    & & \\ \\

    \theutterance \stepcounter{utterance}  
    & & & \multicolumn{2}{p{0.3\linewidth}}{
        \cellcolor[rgb]{0.9,0.9,0.9}{
            \makecell[{{p{\linewidth}}}]{
                \texttt{\tiny{[GM$|$GM]}}
                \texttt{D moved to (7, 1)} \\
            }
        }
    }
    & & \\ \\

    \theutterance \stepcounter{utterance}  
    & & \multicolumn{4}{p{0.6\linewidth}}{
        \cellcolor[rgb]{0.9,0.9,0.9}{
            \makecell[{{p{\linewidth}}}]{
                \texttt{\tiny{[P1$\langle$GM]}}
                \texttt{Edukalt teisaldatud 'P' asukohta (1,2). Teie uuendatud ruudustik näeb välja selline:} \\
\texttt{```} \\
\texttt{ 1234567} \\
\texttt{╔═╤═╤═══╗} \\
\texttt{║S│◌│◌◌◌║ 1} \\
\texttt{║P├─┤T◌D║ 2} \\
\texttt{╟─┤◌├─┬─╢ 3} \\
\texttt{║◌│M│O│A║ 4} \\
\texttt{╟─┼─┘◌│◌║ 5} \\
\texttt{║◌│◌◌┌┴─╢ 6} \\
\texttt{║◌│◌◌│◌◌║ 7} \\
\texttt{╚═╧══╧══╝} \\
\texttt{```} \\
\texttt{Sa mängid praegu vooru 6 maksimaalsest 28 voorust.} \\
\texttt{Teil on kollektiivselt kogunenud 5 karistust 16 karistusest.} \\
\texttt{Teine mängija liigutas objekti oma ruudustikul.} \\
\\ 
\\ 
\texttt{Mis on sinu järgmine käsk?} \\
            }
        }
    }
    & \\ \\

    \theutterance \stepcounter{utterance}  
    & \multicolumn{4}{p{0.6\linewidth}}{
        \cellcolor[rgb]{0.9,0.9,0.9}{
            \makecell[{{p{\linewidth}}}]{
                \texttt{\tiny{[P1$\rangle$GM]}}
                \texttt{LIIGUTA: S, (7, 2)} \\
            }
        }
    }
    & & \\ \\

    \theutterance \stepcounter{utterance}  
    & & & \multicolumn{2}{p{0.3\linewidth}}{
        \cellcolor[rgb]{0.9,0.9,0.9}{
            \makecell[{{p{\linewidth}}}]{
                \texttt{\tiny{[GM$|$GM]}}
                \texttt{LIIGUTA: S, (7, 2)} \\
            }
        }
    }
    & & \\ \\

    \theutterance \stepcounter{utterance}  
    & & & \multicolumn{2}{p{0.3\linewidth}}{
        \cellcolor[rgb]{0.9,0.9,0.9}{
            \makecell[{{p{\linewidth}}}]{
                \texttt{\tiny{[GM$|$GM]}}
                \texttt{Karistus: (7,2) ei ole tühi, vaid sisaldab objekti 'D'.} \\
\\ 
\texttt{Teil on kollektiivselt kogunenud 6 karistust 16 karistusest.} \\
\texttt{Palun proovi uuesti!} \\
            }
        }
    }
    & & \\ \\

    \theutterance \stepcounter{utterance}  
    & & \multicolumn{4}{p{0.6\linewidth}}{
        \cellcolor[rgb]{0.9,0.9,0.9}{
            \makecell[{{p{\linewidth}}}]{
                \texttt{\tiny{[P1$\langle$GM]}}
                \texttt{Karistus: (7,2) ei ole tühi, vaid sisaldab objekti 'D'.} \\
\\ 
\texttt{Teil on kollektiivselt kogunenud 6 karistust 16 karistusest.} \\
\texttt{Palun proovi uuesti!} \\
            }
        }
    }
    & \\ \\

    \theutterance \stepcounter{utterance}  
    & \multicolumn{4}{p{0.6\linewidth}}{
        \cellcolor[rgb]{0.9,0.9,0.9}{
            \makecell[{{p{\linewidth}}}]{
                \texttt{\tiny{[P1$\rangle$GM]}}
                \texttt{LIIGUTA: D, (7, 1)} \\
            }
        }
    }
    & & \\ \\

    \theutterance \stepcounter{utterance}  
    & & & \multicolumn{2}{p{0.3\linewidth}}{
        \cellcolor[rgb]{0.9,0.9,0.9}{
            \makecell[{{p{\linewidth}}}]{
                \texttt{\tiny{[GM$|$GM]}}
                \texttt{LIIGUTA: D, (7, 1)} \\
            }
        }
    }
    & & \\ \\

    \theutterance \stepcounter{utterance}  
    & & & \multicolumn{2}{p{0.3\linewidth}}{
        \cellcolor[rgb]{0.9,0.9,0.9}{
            \makecell[{{p{\linewidth}}}]{
                \texttt{\tiny{[GM$|$GM]}}
                \texttt{D moved to (7, 1)} \\
            }
        }
    }
    & & \\ \\

    \theutterance \stepcounter{utterance}  
    & & \multicolumn{4}{p{0.6\linewidth}}{
        \cellcolor[rgb]{0.9,0.9,0.9}{
            \makecell[{{p{\linewidth}}}]{
                \texttt{\tiny{[GM$\rangle$P2]}}
                \texttt{Edukalt teisaldatud 'D' asukohta (7,1). Teie uuendatud ruudustik näeb välja selline:} \\
\texttt{```} \\
\texttt{ 1234567} \\
\texttt{╔═╤═╤═══╗} \\
\texttt{║P│O│◌◌D║ 1} \\
\texttt{║◌├─┤◌◌◌║ 2} \\
\texttt{╟─┤◌├─┬─╢ 3} \\
\texttt{║◌│◌│◌│A║ 4} \\
\texttt{╟─┼─┘◌│◌║ 5} \\
\texttt{║T│◌◌┌┴─╢ 6} \\
\texttt{║M│◌◌│◌S║ 7} \\
\texttt{╚═╧══╧══╝} \\
\texttt{```} \\
\texttt{Sa mängid praegu vooru 7 maksimaalsest 28 voorust.} \\
\texttt{Teil on kollektiivselt kogunenud 6 karistust 16 karistusest.} \\
\texttt{Teine mängija liigutas objekti oma ruudustikul.} \\
\\ 
\\ 
\texttt{Mis on sinu järgmine käsk?} \\
            }
        }
    }
    & \\ \\

    \theutterance \stepcounter{utterance}  
    & & & \multicolumn{4}{p{0.6\linewidth}}{
        \cellcolor[rgb]{0.9,0.9,0.9}{
            \makecell[{{p{\linewidth}}}]{
                \texttt{\tiny{[GM$\langle$P2]}}
                \texttt{LIIGUTA: A, (6, 1)} \\
            }
        }
    }
     \\ \\

    \theutterance \stepcounter{utterance}  
    & & & \multicolumn{2}{p{0.3\linewidth}}{
        \cellcolor[rgb]{0.9,0.9,0.9}{
            \makecell[{{p{\linewidth}}}]{
                \texttt{\tiny{[GM$|$GM]}}
                \texttt{LIIGUTA: A, (6, 1)} \\
            }
        }
    }
    & & \\ \\

    \theutterance \stepcounter{utterance}  
    & & & \multicolumn{2}{p{0.3\linewidth}}{
        \cellcolor[rgb]{0.9,0.9,0.9}{
            \makecell[{{p{\linewidth}}}]{
                \texttt{\tiny{[GM$|$GM]}}
                \texttt{A moved to (6, 1)} \\
            }
        }
    }
    & & \\ \\

    \theutterance \stepcounter{utterance}  
    & & \multicolumn{4}{p{0.6\linewidth}}{
        \cellcolor[rgb]{0.9,0.9,0.9}{
            \makecell[{{p{\linewidth}}}]{
                \texttt{\tiny{[P1$\langle$GM]}}
                \texttt{Edukalt teisaldatud 'D' asukohta (7,1). Teie uuendatud ruudustik näeb välja selline:} \\
\texttt{```} \\
\texttt{ 1234567} \\
\texttt{╔═╤═╤═══╗} \\
\texttt{║S│◌│◌◌D║ 1} \\
\texttt{║P├─┤T◌◌║ 2} \\
\texttt{╟─┤◌├─┬─╢ 3} \\
\texttt{║◌│M│O│A║ 4} \\
\texttt{╟─┼─┘◌│◌║ 5} \\
\texttt{║◌│◌◌┌┴─╢ 6} \\
\texttt{║◌│◌◌│◌◌║ 7} \\
\texttt{╚═╧══╧══╝} \\
\texttt{```} \\
\texttt{Sa mängid praegu vooru 7 maksimaalsest 28 voorust.} \\
\texttt{Teil on kollektiivselt kogunenud 6 karistust 16 karistusest.} \\
\texttt{Teine mängija liigutas objekti oma ruudustikul.} \\
\\ 
\\ 
\texttt{Mis on sinu järgmine käsk?} \\
            }
        }
    }
    & \\ \\

    \theutterance \stepcounter{utterance}  
    & \multicolumn{4}{p{0.6\linewidth}}{
        \cellcolor[rgb]{0.9,0.9,0.9}{
            \makecell[{{p{\linewidth}}}]{
                \texttt{\tiny{[P1$\rangle$GM]}}
                \texttt{LIIGUTA: S, (7, 2)} \\
            }
        }
    }
    & & \\ \\

    \theutterance \stepcounter{utterance}  
    & & & \multicolumn{2}{p{0.3\linewidth}}{
        \cellcolor[rgb]{0.9,0.9,0.9}{
            \makecell[{{p{\linewidth}}}]{
                \texttt{\tiny{[GM$|$GM]}}
                \texttt{LIIGUTA: S, (7, 2)} \\
            }
        }
    }
    & & \\ \\

    \theutterance \stepcounter{utterance}  
    & & & \multicolumn{2}{p{0.3\linewidth}}{
        \cellcolor[rgb]{0.9,0.9,0.9}{
            \makecell[{{p{\linewidth}}}]{
                \texttt{\tiny{[GM$|$GM]}}
                \texttt{S moved to (7, 2)} \\
            }
        }
    }
    & & \\ \\

    \theutterance \stepcounter{utterance}  
    & & \multicolumn{4}{p{0.6\linewidth}}{
        \cellcolor[rgb]{0.9,0.9,0.9}{
            \makecell[{{p{\linewidth}}}]{
                \texttt{\tiny{[GM$\rangle$P2]}}
                \texttt{Edukalt teisaldatud 'A' asukohta (6,1). Teie uuendatud ruudustik näeb välja selline:} \\
\texttt{```} \\
\texttt{ 1234567} \\
\texttt{╔═╤═╤═══╗} \\
\texttt{║P│O│◌AD║ 1} \\
\texttt{║◌├─┤◌◌◌║ 2} \\
\texttt{╟─┤◌├─┬─╢ 3} \\
\texttt{║◌│◌│◌│◌║ 4} \\
\texttt{╟─┼─┘◌│◌║ 5} \\
\texttt{║T│◌◌┌┴─╢ 6} \\
\texttt{║M│◌◌│◌S║ 7} \\
\texttt{╚═╧══╧══╝} \\
\texttt{```} \\
\texttt{Sa mängid praegu vooru 8 maksimaalsest 28 voorust.} \\
\texttt{Teil on kollektiivselt kogunenud 6 karistust 16 karistusest.} \\
\texttt{Teine mängija liigutas objekti oma ruudustikul.} \\
\\ 
\\ 
\texttt{Mis on sinu järgmine käsk?} \\
            }
        }
    }
    & \\ \\

    \theutterance \stepcounter{utterance}  
    & & & \multicolumn{4}{p{0.6\linewidth}}{
        \cellcolor[rgb]{0.9,0.9,0.9}{
            \makecell[{{p{\linewidth}}}]{
                \texttt{\tiny{[GM$\langle$P2]}}
                \texttt{LIIGUTA: M, (5, 1)} \\
            }
        }
    }
     \\ \\

    \theutterance \stepcounter{utterance}  
    & & & \multicolumn{2}{p{0.3\linewidth}}{
        \cellcolor[rgb]{0.9,0.9,0.9}{
            \makecell[{{p{\linewidth}}}]{
                \texttt{\tiny{[GM$|$GM]}}
                \texttt{LIIGUTA: M, (5, 1)} \\
            }
        }
    }
    & & \\ \\

    \theutterance \stepcounter{utterance}  
    & & & \multicolumn{2}{p{0.3\linewidth}}{
        \cellcolor[rgb]{0.9,0.9,0.9}{
            \makecell[{{p{\linewidth}}}]{
                \texttt{\tiny{[GM$|$GM]}}
                \texttt{M moved to (5, 1)} \\
            }
        }
    }
    & & \\ \\

    \theutterance \stepcounter{utterance}  
    & & \multicolumn{4}{p{0.6\linewidth}}{
        \cellcolor[rgb]{0.9,0.9,0.9}{
            \makecell[{{p{\linewidth}}}]{
                \texttt{\tiny{[P1$\langle$GM]}}
                \texttt{Edukalt teisaldatud 'S' asukohta (7,2). Teie uuendatud ruudustik näeb välja selline:} \\
\texttt{```} \\
\texttt{ 1234567} \\
\texttt{╔═╤═╤═══╗} \\
\texttt{║◌│◌│◌◌D║ 1} \\
\texttt{║P├─┤T◌S║ 2} \\
\texttt{╟─┤◌├─┬─╢ 3} \\
\texttt{║◌│M│O│A║ 4} \\
\texttt{╟─┼─┘◌│◌║ 5} \\
\texttt{║◌│◌◌┌┴─╢ 6} \\
\texttt{║◌│◌◌│◌◌║ 7} \\
\texttt{╚═╧══╧══╝} \\
\texttt{```} \\
\texttt{Sa mängid praegu vooru 8 maksimaalsest 28 voorust.} \\
\texttt{Teil on kollektiivselt kogunenud 6 karistust 16 karistusest.} \\
\texttt{Teine mängija liigutas objekti oma ruudustikul.} \\
\\ 
\\ 
\texttt{Mis on sinu järgmine käsk?} \\
            }
        }
    }
    & \\ \\

    \theutterance \stepcounter{utterance}  
    & \multicolumn{4}{p{0.6\linewidth}}{
        \cellcolor[rgb]{0.9,0.9,0.9}{
            \makecell[{{p{\linewidth}}}]{
                \texttt{\tiny{[P1$\rangle$GM]}}
                \texttt{LIIGUTA: T, (6, 2)} \\
            }
        }
    }
    & & \\ \\

    \theutterance \stepcounter{utterance}  
    & & & \multicolumn{2}{p{0.3\linewidth}}{
        \cellcolor[rgb]{0.9,0.9,0.9}{
            \makecell[{{p{\linewidth}}}]{
                \texttt{\tiny{[GM$|$GM]}}
                \texttt{LIIGUTA: T, (6, 2)} \\
            }
        }
    }
    & & \\ \\

    \theutterance \stepcounter{utterance}  
    & & & \multicolumn{2}{p{0.3\linewidth}}{
        \cellcolor[rgb]{0.9,0.9,0.9}{
            \makecell[{{p{\linewidth}}}]{
                \texttt{\tiny{[GM$|$GM]}}
                \texttt{T moved to (6, 2)} \\
            }
        }
    }
    & & \\ \\

    \theutterance \stepcounter{utterance}  
    & & \multicolumn{4}{p{0.6\linewidth}}{
        \cellcolor[rgb]{0.9,0.9,0.9}{
            \makecell[{{p{\linewidth}}}]{
                \texttt{\tiny{[GM$\rangle$P2]}}
                \texttt{Edukalt teisaldatud 'M' asukohta (5,1). Teie uuendatud ruudustik näeb välja selline:} \\
\texttt{```} \\
\texttt{ 1234567} \\
\texttt{╔═╤═╤═══╗} \\
\texttt{║P│O│MAD║ 1} \\
\texttt{║◌├─┤◌◌◌║ 2} \\
\texttt{╟─┤◌├─┬─╢ 3} \\
\texttt{║◌│◌│◌│◌║ 4} \\
\texttt{╟─┼─┘◌│◌║ 5} \\
\texttt{║T│◌◌┌┴─╢ 6} \\
\texttt{║◌│◌◌│◌S║ 7} \\
\texttt{╚═╧══╧══╝} \\
\texttt{```} \\
\texttt{Sa mängid praegu vooru 9 maksimaalsest 28 voorust.} \\
\texttt{Teil on kollektiivselt kogunenud 6 karistust 16 karistusest.} \\
\texttt{Teine mängija liigutas objekti oma ruudustikul.} \\
\\ 
\\ 
\texttt{Mis on sinu järgmine käsk?} \\
            }
        }
    }
    & \\ \\

    \theutterance \stepcounter{utterance}  
    & & & \multicolumn{4}{p{0.6\linewidth}}{
        \cellcolor[rgb]{0.9,0.9,0.9}{
            \makecell[{{p{\linewidth}}}]{
                \texttt{\tiny{[GM$\langle$P2]}}
                \texttt{LIIGUTA: P, (1, 2)} \\
            }
        }
    }
     \\ \\

    \theutterance \stepcounter{utterance}  
    & & & \multicolumn{2}{p{0.3\linewidth}}{
        \cellcolor[rgb]{0.9,0.9,0.9}{
            \makecell[{{p{\linewidth}}}]{
                \texttt{\tiny{[GM$|$GM]}}
                \texttt{LIIGUTA: P, (1, 2)} \\
            }
        }
    }
    & & \\ \\

    \theutterance \stepcounter{utterance}  
    & & & \multicolumn{2}{p{0.3\linewidth}}{
        \cellcolor[rgb]{0.9,0.9,0.9}{
            \makecell[{{p{\linewidth}}}]{
                \texttt{\tiny{[GM$|$GM]}}
                \texttt{P moved to (1, 2)} \\
            }
        }
    }
    & & \\ \\

    \theutterance \stepcounter{utterance}  
    & & \multicolumn{4}{p{0.6\linewidth}}{
        \cellcolor[rgb]{0.9,0.9,0.9}{
            \makecell[{{p{\linewidth}}}]{
                \texttt{\tiny{[P1$\langle$GM]}}
                \texttt{Edukalt teisaldatud 'T' asukohta (6,2). Teie uuendatud ruudustik näeb välja selline:} \\
\texttt{```} \\
\texttt{ 1234567} \\
\texttt{╔═╤═╤═══╗} \\
\texttt{║◌│◌│◌◌D║ 1} \\
\texttt{║P├─┤◌TS║ 2} \\
\texttt{╟─┤◌├─┬─╢ 3} \\
\texttt{║◌│M│O│A║ 4} \\
\texttt{╟─┼─┘◌│◌║ 5} \\
\texttt{║◌│◌◌┌┴─╢ 6} \\
\texttt{║◌│◌◌│◌◌║ 7} \\
\texttt{╚═╧══╧══╝} \\
\texttt{```} \\
\texttt{Sa mängid praegu vooru 9 maksimaalsest 28 voorust.} \\
\texttt{Teil on kollektiivselt kogunenud 6 karistust 16 karistusest.} \\
\texttt{Teine mängija liigutas objekti oma ruudustikul.} \\
\\ 
\\ 
\texttt{Mis on sinu järgmine käsk?} \\
            }
        }
    }
    & \\ \\

    \theutterance \stepcounter{utterance}  
    & \multicolumn{4}{p{0.6\linewidth}}{
        \cellcolor[rgb]{0.9,0.9,0.9}{
            \makecell[{{p{\linewidth}}}]{
                \texttt{\tiny{[P1$\rangle$GM]}}
                \texttt{LIIGUTA: M, (5, 1)} \\
            }
        }
    }
    & & \\ \\

    \theutterance \stepcounter{utterance}  
    & & & \multicolumn{2}{p{0.3\linewidth}}{
        \cellcolor[rgb]{0.9,0.9,0.9}{
            \makecell[{{p{\linewidth}}}]{
                \texttt{\tiny{[GM$|$GM]}}
                \texttt{LIIGUTA: M, (5, 1)} \\
            }
        }
    }
    & & \\ \\

    \theutterance \stepcounter{utterance}  
    & & & \multicolumn{2}{p{0.3\linewidth}}{
        \cellcolor[rgb]{0.9,0.9,0.9}{
            \makecell[{{p{\linewidth}}}]{
                \texttt{\tiny{[GM$|$GM]}}
                \texttt{M moved to (5, 1)} \\
            }
        }
    }
    & & \\ \\

    \theutterance \stepcounter{utterance}  
    & & \multicolumn{4}{p{0.6\linewidth}}{
        \cellcolor[rgb]{0.9,0.9,0.9}{
            \makecell[{{p{\linewidth}}}]{
                \texttt{\tiny{[GM$\rangle$P2]}}
                \texttt{Edukalt teisaldatud 'P' asukohta (1,2). Teie uuendatud ruudustik näeb välja selline:} \\
\texttt{```} \\
\texttt{ 1234567} \\
\texttt{╔═╤═╤═══╗} \\
\texttt{║◌│O│MAD║ 1} \\
\texttt{║P├─┤◌◌◌║ 2} \\
\texttt{╟─┤◌├─┬─╢ 3} \\
\texttt{║◌│◌│◌│◌║ 4} \\
\texttt{╟─┼─┘◌│◌║ 5} \\
\texttt{║T│◌◌┌┴─╢ 6} \\
\texttt{║◌│◌◌│◌S║ 7} \\
\texttt{╚═╧══╧══╝} \\
\texttt{```} \\
\texttt{Sa mängid praegu vooru 10 maksimaalsest 28 voorust.} \\
\texttt{Teil on kollektiivselt kogunenud 6 karistust 16 karistusest.} \\
\texttt{Teine mängija liigutas objekti oma ruudustikul.} \\
\\ 
\\ 
\texttt{Mis on sinu järgmine käsk?} \\
            }
        }
    }
    & \\ \\

    \theutterance \stepcounter{utterance}  
    & & & \multicolumn{4}{p{0.6\linewidth}}{
        \cellcolor[rgb]{0.9,0.9,0.9}{
            \makecell[{{p{\linewidth}}}]{
                \texttt{\tiny{[GM$\langle$P2]}}
                \texttt{LIIGUTA: O, (5, 2)} \\
            }
        }
    }
     \\ \\

    \theutterance \stepcounter{utterance}  
    & & & \multicolumn{2}{p{0.3\linewidth}}{
        \cellcolor[rgb]{0.9,0.9,0.9}{
            \makecell[{{p{\linewidth}}}]{
                \texttt{\tiny{[GM$|$GM]}}
                \texttt{LIIGUTA: O, (5, 2)} \\
            }
        }
    }
    & & \\ \\

    \theutterance \stepcounter{utterance}  
    & & & \multicolumn{2}{p{0.3\linewidth}}{
        \cellcolor[rgb]{0.9,0.9,0.9}{
            \makecell[{{p{\linewidth}}}]{
                \texttt{\tiny{[GM$|$GM]}}
                \texttt{O moved to (5, 2)} \\
            }
        }
    }
    & & \\ \\

    \theutterance \stepcounter{utterance}  
    & & \multicolumn{4}{p{0.6\linewidth}}{
        \cellcolor[rgb]{0.9,0.9,0.9}{
            \makecell[{{p{\linewidth}}}]{
                \texttt{\tiny{[P1$\langle$GM]}}
                \texttt{Edukalt teisaldatud 'M' asukohta (5,1). Teie uuendatud ruudustik näeb välja selline:} \\
\texttt{```} \\
\texttt{ 1234567} \\
\texttt{╔═╤═╤═══╗} \\
\texttt{║◌│◌│M◌D║ 1} \\
\texttt{║P├─┤◌TS║ 2} \\
\texttt{╟─┤◌├─┬─╢ 3} \\
\texttt{║◌│◌│O│A║ 4} \\
\texttt{╟─┼─┘◌│◌║ 5} \\
\texttt{║◌│◌◌┌┴─╢ 6} \\
\texttt{║◌│◌◌│◌◌║ 7} \\
\texttt{╚═╧══╧══╝} \\
\texttt{```} \\
\texttt{Sa mängid praegu vooru 10 maksimaalsest 28 voorust.} \\
\texttt{Teil on kollektiivselt kogunenud 6 karistust 16 karistusest.} \\
\texttt{Teine mängija liigutas objekti oma ruudustikul.} \\
\\ 
\\ 
\texttt{Mis on sinu järgmine käsk?} \\
            }
        }
    }
    & \\ \\

    \theutterance \stepcounter{utterance}  
    & \multicolumn{4}{p{0.6\linewidth}}{
        \cellcolor[rgb]{0.9,0.9,0.9}{
            \makecell[{{p{\linewidth}}}]{
                \texttt{\tiny{[P1$\rangle$GM]}}
                \texttt{LIIGUTA: O, (5, 2)} \\
            }
        }
    }
    & & \\ \\

    \theutterance \stepcounter{utterance}  
    & & & \multicolumn{2}{p{0.3\linewidth}}{
        \cellcolor[rgb]{0.9,0.9,0.9}{
            \makecell[{{p{\linewidth}}}]{
                \texttt{\tiny{[GM$|$GM]}}
                \texttt{LIIGUTA: O, (5, 2)} \\
            }
        }
    }
    & & \\ \\

    \theutterance \stepcounter{utterance}  
    & & & \multicolumn{2}{p{0.3\linewidth}}{
        \cellcolor[rgb]{0.9,0.9,0.9}{
            \makecell[{{p{\linewidth}}}]{
                \texttt{\tiny{[GM$|$GM]}}
                \texttt{O moved to (5, 2)} \\
            }
        }
    }
    & & \\ \\

    \theutterance \stepcounter{utterance}  
    & & \multicolumn{4}{p{0.6\linewidth}}{
        \cellcolor[rgb]{0.9,0.9,0.9}{
            \makecell[{{p{\linewidth}}}]{
                \texttt{\tiny{[GM$\rangle$P2]}}
                \texttt{Edukalt teisaldatud 'O' asukohta (5,2). Teie uuendatud ruudustik näeb välja selline:} \\
\texttt{```} \\
\texttt{ 1234567} \\
\texttt{╔═╤═╤═══╗} \\
\texttt{║◌│◌│MAD║ 1} \\
\texttt{║P├─┤O◌◌║ 2} \\
\texttt{╟─┤◌├─┬─╢ 3} \\
\texttt{║◌│◌│◌│◌║ 4} \\
\texttt{╟─┼─┘◌│◌║ 5} \\
\texttt{║T│◌◌┌┴─╢ 6} \\
\texttt{║◌│◌◌│◌S║ 7} \\
\texttt{╚═╧══╧══╝} \\
\texttt{```} \\
\texttt{Sa mängid praegu vooru 11 maksimaalsest 28 voorust.} \\
\texttt{Teil on kollektiivselt kogunenud 6 karistust 16 karistusest.} \\
\texttt{Teine mängija liigutas objekti oma ruudustikul.} \\
\\ 
\\ 
\texttt{Mis on sinu järgmine käsk?} \\
            }
        }
    }
    & \\ \\

    \theutterance \stepcounter{utterance}  
    & & & \multicolumn{4}{p{0.6\linewidth}}{
        \cellcolor[rgb]{0.9,0.9,0.9}{
            \makecell[{{p{\linewidth}}}]{
                \texttt{\tiny{[GM$\langle$P2]}}
                \texttt{LIIGUTA: T, (6, 2)} \\
            }
        }
    }
     \\ \\

    \theutterance \stepcounter{utterance}  
    & & & \multicolumn{2}{p{0.3\linewidth}}{
        \cellcolor[rgb]{0.9,0.9,0.9}{
            \makecell[{{p{\linewidth}}}]{
                \texttt{\tiny{[GM$|$GM]}}
                \texttt{LIIGUTA: T, (6, 2)} \\
            }
        }
    }
    & & \\ \\

    \theutterance \stepcounter{utterance}  
    & & & \multicolumn{2}{p{0.3\linewidth}}{
        \cellcolor[rgb]{0.9,0.9,0.9}{
            \makecell[{{p{\linewidth}}}]{
                \texttt{\tiny{[GM$|$GM]}}
                \texttt{T moved to (6, 2)} \\
            }
        }
    }
    & & \\ \\

    \theutterance \stepcounter{utterance}  
    & & \multicolumn{4}{p{0.6\linewidth}}{
        \cellcolor[rgb]{0.9,0.9,0.9}{
            \makecell[{{p{\linewidth}}}]{
                \texttt{\tiny{[P1$\langle$GM]}}
                \texttt{Edukalt teisaldatud 'O' asukohta (5,2). Teie uuendatud ruudustik näeb välja selline:} \\
\texttt{```} \\
\texttt{ 1234567} \\
\texttt{╔═╤═╤═══╗} \\
\texttt{║◌│◌│M◌D║ 1} \\
\texttt{║P├─┤OTS║ 2} \\
\texttt{╟─┤◌├─┬─╢ 3} \\
\texttt{║◌│◌│◌│A║ 4} \\
\texttt{╟─┼─┘◌│◌║ 5} \\
\texttt{║◌│◌◌┌┴─╢ 6} \\
\texttt{║◌│◌◌│◌◌║ 7} \\
\texttt{╚═╧══╧══╝} \\
\texttt{```} \\
\texttt{Sa mängid praegu vooru 11 maksimaalsest 28 voorust.} \\
\texttt{Teil on kollektiivselt kogunenud 6 karistust 16 karistusest.} \\
\texttt{Teine mängija liigutas objekti oma ruudustikul.} \\
\\ 
\\ 
\texttt{Mis on sinu järgmine käsk?} \\
            }
        }
    }
    & \\ \\

    \theutterance \stepcounter{utterance}  
    & \multicolumn{4}{p{0.6\linewidth}}{
        \cellcolor[rgb]{0.9,0.9,0.9}{
            \makecell[{{p{\linewidth}}}]{
                \texttt{\tiny{[P1$\rangle$GM]}}
                \texttt{LIIGUTA: A, (6, 1)} \\
            }
        }
    }
    & & \\ \\

    \theutterance \stepcounter{utterance}  
    & & & \multicolumn{2}{p{0.3\linewidth}}{
        \cellcolor[rgb]{0.9,0.9,0.9}{
            \makecell[{{p{\linewidth}}}]{
                \texttt{\tiny{[GM$|$GM]}}
                \texttt{LIIGUTA: A, (6, 1)} \\
            }
        }
    }
    & & \\ \\

    \theutterance \stepcounter{utterance}  
    & & & \multicolumn{2}{p{0.3\linewidth}}{
        \cellcolor[rgb]{0.9,0.9,0.9}{
            \makecell[{{p{\linewidth}}}]{
                \texttt{\tiny{[GM$|$GM]}}
                \texttt{A moved to (6, 1)} \\
            }
        }
    }
    & & \\ \\

    \theutterance \stepcounter{utterance}  
    & & \multicolumn{4}{p{0.6\linewidth}}{
        \cellcolor[rgb]{0.9,0.9,0.9}{
            \makecell[{{p{\linewidth}}}]{
                \texttt{\tiny{[GM$\rangle$P2]}}
                \texttt{Edukalt teisaldatud 'T' asukohta (6,2). Teie uuendatud ruudustik näeb välja selline:} \\
\texttt{```} \\
\texttt{ 1234567} \\
\texttt{╔═╤═╤═══╗} \\
\texttt{║◌│◌│MAD║ 1} \\
\texttt{║P├─┤OT◌║ 2} \\
\texttt{╟─┤◌├─┬─╢ 3} \\
\texttt{║◌│◌│◌│◌║ 4} \\
\texttt{╟─┼─┘◌│◌║ 5} \\
\texttt{║◌│◌◌┌┴─╢ 6} \\
\texttt{║◌│◌◌│◌S║ 7} \\
\texttt{╚═╧══╧══╝} \\
\texttt{```} \\
\texttt{Sa mängid praegu vooru 12 maksimaalsest 28 voorust.} \\
\texttt{Teil on kollektiivselt kogunenud 6 karistust 16 karistusest.} \\
\texttt{Teine mängija liigutas objekti oma ruudustikul.} \\
\\ 
\\ 
\texttt{Mis on sinu järgmine käsk?} \\
            }
        }
    }
    & \\ \\

    \theutterance \stepcounter{utterance}  
    & & & \multicolumn{4}{p{0.6\linewidth}}{
        \cellcolor[rgb]{0.9,0.9,0.9}{
            \makecell[{{p{\linewidth}}}]{
                \texttt{\tiny{[GM$\langle$P2]}}
                \texttt{LIIGUTA: S, (7, 2)} \\
            }
        }
    }
     \\ \\

    \theutterance \stepcounter{utterance}  
    & & & \multicolumn{2}{p{0.3\linewidth}}{
        \cellcolor[rgb]{0.9,0.9,0.9}{
            \makecell[{{p{\linewidth}}}]{
                \texttt{\tiny{[GM$|$GM]}}
                \texttt{LIIGUTA: S, (7, 2)} \\
            }
        }
    }
    & & \\ \\

    \theutterance \stepcounter{utterance}  
    & & & \multicolumn{2}{p{0.3\linewidth}}{
        \cellcolor[rgb]{0.9,0.9,0.9}{
            \makecell[{{p{\linewidth}}}]{
                \texttt{\tiny{[GM$|$GM]}}
                \texttt{S moved to (7, 2)} \\
            }
        }
    }
    & & \\ \\

    \theutterance \stepcounter{utterance}  
    & & \multicolumn{4}{p{0.6\linewidth}}{
        \cellcolor[rgb]{0.9,0.9,0.9}{
            \makecell[{{p{\linewidth}}}]{
                \texttt{\tiny{[P1$\langle$GM]}}
                \texttt{Edukalt teisaldatud 'A' asukohta (6,1). Teie uuendatud ruudustik näeb välja selline:} \\
\texttt{```} \\
\texttt{ 1234567} \\
\texttt{╔═╤═╤═══╗} \\
\texttt{║◌│◌│MAD║ 1} \\
\texttt{║P├─┤OTS║ 2} \\
\texttt{╟─┤◌├─┬─╢ 3} \\
\texttt{║◌│◌│◌│◌║ 4} \\
\texttt{╟─┼─┘◌│◌║ 5} \\
\texttt{║◌│◌◌┌┴─╢ 6} \\
\texttt{║◌│◌◌│◌◌║ 7} \\
\texttt{╚═╧══╧══╝} \\
\texttt{```} \\
\texttt{Sa mängid praegu vooru 12 maksimaalsest 28 voorust.} \\
\texttt{Teil on kollektiivselt kogunenud 6 karistust 16 karistusest.} \\
\texttt{Teine mängija liigutas objekti oma ruudustikul.} \\
\\ 
\\ 
\texttt{Mis on sinu järgmine käsk?} \\
            }
        }
    }
    & \\ \\

    \theutterance \stepcounter{utterance}  
    & \multicolumn{4}{p{0.6\linewidth}}{
        \cellcolor[rgb]{0.9,0.9,0.9}{
            \makecell[{{p{\linewidth}}}]{
                \texttt{\tiny{[P1$\rangle$GM]}}
                \texttt{ÜTLE: Kõik objektid on paigutatud: P(1,2), O(5,2), T(6,2), S(7,2), D(7,1), A(6,1), M(5,1). Kas oled sa ka valmis? valmis?} \\
            }
        }
    }
    & & \\ \\

    \theutterance \stepcounter{utterance}  
    & & & \multicolumn{2}{p{0.3\linewidth}}{
        \cellcolor[rgb]{0.9,0.9,0.9}{
            \makecell[{{p{\linewidth}}}]{
                \texttt{\tiny{[GM$|$GM]}}
                \texttt{ÜTLE: Kõik objektid on paigutatud: P(1,2), O(5,2), T(6,2), S(7,2), D(7,1), A(6,1), M(5,1). Kas oled sa ka valmis? valmis?} \\
            }
        }
    }
    & & \\ \\

    \theutterance \stepcounter{utterance}  
    & & \multicolumn{4}{p{0.6\linewidth}}{
        \cellcolor[rgb]{0.9,0.9,0.9}{
            \makecell[{{p{\linewidth}}}]{
                \texttt{\tiny{[GM$\rangle$P2]}}
                \texttt{Edukalt teisaldatud 'S' asukohta (7,2). Teie uuendatud ruudustik näeb välja selline:} \\
\texttt{```} \\
\texttt{ 1234567} \\
\texttt{╔═╤═╤═══╗} \\
\texttt{║◌│◌│MAD║ 1} \\
\texttt{║P├─┤OTS║ 2} \\
\texttt{╟─┤◌├─┬─╢ 3} \\
\texttt{║◌│◌│◌│◌║ 4} \\
\texttt{╟─┼─┘◌│◌║ 5} \\
\texttt{║◌│◌◌┌┴─╢ 6} \\
\texttt{║◌│◌◌│◌◌║ 7} \\
\texttt{╚═╧══╧══╝} \\
\texttt{```} \\
\texttt{Sa mängid praegu vooru 13 maksimaalsest 28 voorust.} \\
\texttt{Teil on kollektiivselt kogunenud 6 karistust 16 karistusest.} \\
\texttt{Teine mängija saatis selle sõnumi:} \\
\\ 
\\ 
\texttt{"Kõik objektid on paigutatud: P(1,2), O(5,2), T(6,2), S(7,2), D(7,1), A(6,1), M(5,1). Kas oled sa ka valmis? valmis?"} \\
\\ 
\\ 
\texttt{Mis on sinu järgmine käsk?} \\
            }
        }
    }
    & \\ \\

    \theutterance \stepcounter{utterance}  
    & & & \multicolumn{4}{p{0.6\linewidth}}{
        \cellcolor[rgb]{0.9,0.9,0.9}{
            \makecell[{{p{\linewidth}}}]{
                \texttt{\tiny{[GM$\langle$P2]}}
                \texttt{ÜTLE: valmis!} \\
            }
        }
    }
     \\ \\

    \theutterance \stepcounter{utterance}  
    & & & \multicolumn{2}{p{0.3\linewidth}}{
        \cellcolor[rgb]{0.9,0.9,0.9}{
            \makecell[{{p{\linewidth}}}]{
                \texttt{\tiny{[GM$|$GM]}}
                \texttt{ÜTLE: valmis!} \\
            }
        }
    }
    & & \\ \\

    \theutterance \stepcounter{utterance}  
    & & & \multicolumn{2}{p{0.3\linewidth}}{
        \cellcolor[rgb]{0.9,0.9,0.9}{
            \makecell[{{p{\linewidth}}}]{
                \texttt{\tiny{[GM$|$GM]}}
                \texttt{true} \\
            }
        }
    }
    & & \\ \\

    \theutterance \stepcounter{utterance}  
    & & & \multicolumn{2}{p{0.3\linewidth}}{
        \cellcolor[rgb]{0.9,0.9,0.9}{
            \makecell[{{p{\linewidth}}}]{
                \texttt{\tiny{[GM$|$GM]}}
                \texttt{Initial states:} \\
\texttt{```} \\
\texttt{Player 1:} \\
\texttt{ 1234567} \\
\texttt{╔═╤═╤═══╗} \\
\texttt{║◌│◌│◌◌◌║ 1} \\
\texttt{║S├─┤T◌D║ 2} \\
\texttt{╟─┤◌├─┬─╢ 3} \\
\texttt{║◌│M│O│A║ 4} \\
\texttt{╟─┼─┘P│◌║ 5} \\
\texttt{║◌│◌◌┌┴─╢ 6} \\
\texttt{║◌│◌◌│◌◌║ 7} \\
\texttt{╚═╧══╧══╝} \\
\\ 
\\ 
\texttt{Player 2:} \\
\texttt{ 1234567} \\
\texttt{╔═╤═╤═══╗} \\
\texttt{║◌│◌│◌PO║ 1} \\
\texttt{║◌├─┤◌◌◌║ 2} \\
\texttt{╟─┤◌├─┬─╢ 3} \\
\texttt{║◌│◌│◌│A║ 4} \\
\texttt{╟─┼─┘D│◌║ 5} \\
\texttt{║T│◌◌┌┴─╢ 6} \\
\texttt{║M│◌◌│◌S║ 7} \\
\texttt{╚═╧══╧══╝} \\
\\ 
\texttt{```} \\
            }
        }
    }
    & & \\ \\

    \theutterance \stepcounter{utterance}  
    & & & \multicolumn{2}{p{0.3\linewidth}}{
        \cellcolor[rgb]{0.9,0.9,0.9}{
            \makecell[{{p{\linewidth}}}]{
                \texttt{\tiny{[GM$|$GM]}}
                \texttt{End states:} \\
\texttt{```} \\
\texttt{Player 1:} \\
\texttt{ 1234567} \\
\texttt{╔═╤═╤═══╗} \\
\texttt{║◌│◌│MAD║ 1} \\
\texttt{║P├─┤OTS║ 2} \\
\texttt{╟─┤◌├─┬─╢ 3} \\
\texttt{║◌│◌│◌│◌║ 4} \\
\texttt{╟─┼─┘◌│◌║ 5} \\
\texttt{║◌│◌◌┌┴─╢ 6} \\
\texttt{║◌│◌◌│◌◌║ 7} \\
\texttt{╚═╧══╧══╝} \\
\\ 
\\ 
\texttt{Player 2:} \\
\texttt{ 1234567} \\
\texttt{╔═╤═╤═══╗} \\
\texttt{║◌│◌│MAD║ 1} \\
\texttt{║P├─┤OTS║ 2} \\
\texttt{╟─┤◌├─┬─╢ 3} \\
\texttt{║◌│◌│◌│◌║ 4} \\
\texttt{╟─┼─┘◌│◌║ 5} \\
\texttt{║◌│◌◌┌┴─╢ 6} \\
\texttt{║◌│◌◌│◌◌║ 7} \\
\texttt{╚═╧══╧══╝} \\
\\ 
\texttt{```} \\
            }
        }
    }
    & & \\ \\

    \theutterance \stepcounter{utterance}  
    & & & \multicolumn{2}{p{0.3\linewidth}}{
        \cellcolor[rgb]{0.9,0.9,0.9}{
            \makecell[{{p{\linewidth}}}]{
                \texttt{\tiny{[GM$|$GM]}}
                \texttt{* success: True} \\
\texttt{* lose: False} \\
\texttt{* aborted: False} \\
\texttt{{-}{-}{-}{-}{-}{-}{-}} \\
\texttt{* Shifts: 14.00} \\
\texttt{* Max Shifts: 12.00} \\
\texttt{* Min Shifts: 6.00} \\
\texttt{* End Distance Sum: 0.00} \\
\texttt{* Init Distance Sum: 28.41} \\
\texttt{* Expected Distance Sum: 29.33} \\
\texttt{* Penalties: 6.00} \\
\texttt{* Max Penalties: 16.00} \\
\texttt{* Rounds: 13.00} \\
\texttt{* Max Rounds: 28.00} \\
\texttt{* Object Count: 7.00} \\
            }
        }
    }
    & & \\ \\

\end{supertabular}
}

\end{document}
